% !TEX root = ../main.tex

% --- GENERAL INTRODUCTION (Before Chapter 1) ---
\clearpage
\pagenumbering{arabic}
\setcounter{page}{1}
\section*{Introduction}
\phantomsection
\markboth{INTRODUCTION}{}
\addcontentsline{toc}{section}{Introduction}

The recent evolution of Artificial Intelligence has marked a major shift from traditional conversational models toward \textit{agentic systems}. These systems are characterized by their ability to reason, use external tools, interact with data sources, and execute complex multi-step workflows with a certain level of autonomy. In this context, the subject of this internship focuses on the design and implementation of \textbf{VIP -- Value Intelligent Platform}, an enterprise-grade Agentic AI Hub intended to build, deploy, and operate intelligent applications at scale across different business domains.

\vspace{1em}

Despite the significant progress of Large Language Models, their integration into enterprise environments remains challenging. The main difficulty lies in transforming AI prototypes into reliable, reusable, and governable systems that can be deployed in real operational contexts. Traditional approaches often rely on hard-coded agents, making them difficult to adapt, maintain, or extend. The problem addressed in this project is therefore to design a modular platform capable of generating specialized agents from declarative configurations, while ensuring scalability, traceability, and operational control.

\vspace{1em}

The adopted approach is based on a decoupled and containerized multi-agent architecture. Each agent is designed as an independent FastAPI service, deployed in a Docker container, and configured through prompts, tools, model parameters, and input/output schemas. MLflow is used to manage and version configuration artifacts such as prompts and schemas, while an orchestration layer coordinates the execution of agent workflows. A semantic adapter is also introduced to manage schema transformations between agents and ensure smooth communication across the platform.

\vspace{1em}

This report is organized into five chapters. Chapter 1 presents the internship context, the host company, the project overview, the problem statement, the main objectives, the technology choices, and the adopted methodology. Chapter 2 presents the state of the art and project specifications. It introduces AI agents and multi-agent systems, analyzes existing platforms, defines the VIP requirements and success criteria, and explains the motivation behind the SLM-Evals second project. Chapter 3 presents the system design and architecture of the VIP platform, including the Agent Factory, generated agent runtime, orchestration layer, deployment model, data governance, and observability architecture. Chapter 4 details the implementation, evaluation, and obtained results for both VIP and SLM-Evals. Finally, Chapter 5 concludes the report by summarizing the main contributions, identifying current limitations, and proposing perspectives for future work.
